\scanpage{7}
\uline{connected} \uline{oriented} ``comes to the same'' as giving oneself an $n\in\mathbb{N}$, $n\geq2$, and \inserted{such} a torsor under $\mathbb{Z}/n\mathbb{Z}$ (namely $S$, \unsure{or|--} $A$ \unsure{at one's|--} choice -- NB $\unsure{\sigma|r}:S\longrightarrow A$ is compatible with the action of $G^{\unsure{0|\circ}}=\mathbb{Z}/n\mathbb{Z}$ i.e. with the action of $u_{\omega}$.

Giving oneself a \uline{connected contour} is equivalent to : the data of

$1^{\circ})$ a cyclic group $G^{0}$ (of order $n\geq2$)

$2^{\circ})$ a \deleted{\unsure{set|a}} subset of $G^{0}$ \deleted{\unsure{formed} of} of the form $\{g,g^{-1}\}$, where $g$ is a generator \inserted{(\unsure{thus|of the} subset of cardinality $2$ \unsure{if and only if|ss|if and only if} $n\ne2$)}

$3^{\circ})$ A torsor $S$ under $G^{0}$

$\bullet$ At least if $n\geq3$ \deleted{\unsure{otherwise|if one}} ; for a description

$\bullet$ general, one must in addition give oneself, in place of $2^{\circ}$, a

$2$ bis) a set $\Omega$ with two \unsure{elements|elements}, and a map $\Omega\xrightarrow{i}G$ such that, if $\Omega=\{a,b\}$, one has $i(a)i(b)=1$, $i(a)$ generator of $G$

\begin{center}\rule{1cm}{.35pt}\end{center}

Oriented contour (not necessarily connected) $=$ set on which $\mathbb{Z}$ acts (i.e \unsure{with|or} automorphism $\unsure{u|n}$)

Contour (not necessarily connected) $=$ \ldots
% archive sheet number: 6
